\chapter{The Proprietary Market-Making Firm}\label{ch:fm:the-proprietary-market-making-firm}

In July 2024 the board of a listed exchange-traded-product market maker announced that it would stop paying regular dividends. The money was to
stay in the firm as trading capital: the capital its prime brokers require it to keep with them before it may trade, and which sets how much it
can trade. Eighteen months later that capital stood at \euro1\,044 million, a third more than a year earlier, and the firm had borrowed \$200
million on top. A proprietary trading firm has no clients' money and no fees; everything it can trade is what its owners have left in it, and
what it can borrow against that. How much to leave in, and how much to spend each year to stay fast enough to earn anything, are the two decisions
this chapter models.

\section{Ownership: partnerships, private companies and listed firms}

A proprietary trading firm (One Quant Book 1, chapter 1) is owned by the people who founded it and those they have admitted since; a few are
listed. Ownership decides three things: who receives the profit, who can withdraw capital and when, and who bears a loss.

\begin{definition}[Limited liability partnership, members' capital]\label{def:fm:the-proprietary-market-making-firm:llp}
A \emph{limited liability partnership}\index{limited liability partnership} is a body corporate, with legal personality separate from its members,
whose members share its profit under a members' agreement and whose liability is limited to what they have contributed. A member's
\emph{members' capital}\index{members' capital} is the capital the member has contributed or left in the firm, recorded in the member's capital
account and repayable on departure on the terms of the agreement.
\end{definition}

Private companies owned by their founders and staff work the same way in substance: profit is distributed by dividends or by bonuses instead of
allocations, and capital is equity instead of \omterm{def:fm:the-proprietary-market-making-firm:llp}{members' capital}. A listed firm adds outside shareholders who expect a dividend policy and a share
price, which is why its decisions to retain profit are announced, as in the hook, and a private firm's are not.

\begin{omfigure}
\begin{tikzpicture}[font=\footnotesize,>=stealth,box/.style={draw,rounded corners,minimum width=2.6cm,minimum height=0.62cm,align=center,
  fill=omDef!10}]
\node[box] (rev) at (0,3.2) {revenue\\{\scriptsize capture $\times$ volume}};
\node[box] (prof) at (3.6,3.2) {profit\\{\scriptsize after costs and pay}};
\node[box,fill=omProp!12] (alloc) at (7.4,3.2) {allocation\\{\scriptsize by points}};
\node[box,fill=omThm!10] (draw) at (7.4,1.6) {drawings\\{\scriptsize payout}};
\node[box] (cap) at (3.6,1.6) {members' capital\\{\scriptsize retained}};
\node[box] (tc) at (3.6,0) {trading capital\\{\scriptsize with prime brokers}};
\node[box] (vol) at (0,0) {volume\\{\scriptsize capital / margin}};
\node[box,fill=gray!12] (tech) at (0,1.6) {technology\\{\scriptsize keeps capture}};
\draw[->] (rev) -- (prof); \draw[->] (prof) -- (alloc); \draw[->] (alloc) -- (draw);
\draw[->] (alloc.south west) -- (cap.north east); \draw[->] (cap) -- (tc); \draw[->] (tc) -- (vol);
\draw[->] (vol) -- (tech) node[midway,left,font=\scriptsize] {}; \draw[->] (tech) -- (rev);
\draw[->,dashed,gray] (prof.south west) -- node[above,sloped,font=\scriptsize] {spending} (tech.north east);
\end{tikzpicture}

\omcaption{A trading partnership's year. Profit is allocated by points; what members do not draw stays as capital, which the prime brokers turn
into trading capacity; technology spending keeps capture, and so revenue, from decaying. Model: \texttt{firm.partnership}.}\label{fig:fm:the-proprietary-market-making-firm:cycle}
\end{omfigure}

The members' agreement fixes each member's \emph{points}, the share of profit allocated to them, and how much of an allocation may be drawn. A
departing member's capital is typically repaid over several years, so that a departure does not drain the trading capital at once. That delay is
part of the firm's capital planning: a departure is a withdrawal of capital scheduled in advance.

\begin{example}[An allocation]\label{ex:fm:the-proprietary-market-making-firm:alloc}
Four members hold 40, 30, 20 and 10 points and draw 80\% of their allocations. A year's profit of \$116.25 million allocates \$46.5 million to the
first; she draws \$37.2 million and \$9.3 million stays in her capital account. In all the members draw \$93.0 million and leave \$23.25 million,
a \omterm{def:fm:the-proprietary-market-making-firm:retention}{retention ratio} of 20\%.
\end{example}

\begin{definition}[Retention ratio]\label{def:fm:the-proprietary-market-making-firm:retention}
A firm's \emph{retention ratio}\index{retention ratio} is the share of a year's profit it keeps as capital instead of distributing to its owners;
one minus it is the payout ratio.
\end{definition}

\section{Capital: what it is for and how much is required}

A proprietary firm holds capital for three masters: its regulator, its clearing firms and prime brokers, and itself. The regulator's requirement
protects the market from the firm's failure; the clearing firm's and the prime broker's protect them from the firm's losses before they can close
it out; the firm's own buffer keeps it trading through the losses it expects.

\begin{definition}[Own funds, fixed overheads requirement, K-factor requirement]\label{def:fm:the-proprietary-market-making-firm:ownfunds}
An investment firm's \emph{own funds}\index{own funds} are the regulatory capital it counts against its requirements: mostly paid-up equity and
retained earnings. Under the European rules for investment firms its requirement is the highest of a permanent minimum, a \emph{fixed overheads
requirement}\index{fixed overheads requirement}, one quarter of the preceding year's fixed overheads, and a \emph{K-factor
requirement}\index{K-factor requirement}, the sum of charges proportional to measures of its activity (for a dealer on its own account: its daily
trading flow, and the risk of its net positions or the margin it posts to clearing).
\end{definition}

\begin{dated}{2026-09}{The capital rules for proprietary trading firms}\label{dat:fm:the-proprietary-market-making-firm:rules}
\textbf{European Union.} Regulation (EU) 2019/2033 (the Investment Firms Regulation), article 11: \omterm{def:fm:the-proprietary-market-making-firm:ownfunds}{own funds} of at least the highest of the \omterm{def:fm:the-proprietary-market-making-firm:ownfunds}{fixed
overheads requirement} (article 13: one quarter of the preceding year's fixed overheads), the permanent minimum (article 14, the initial capital
of Directive (EU) 2019/2034, article 9: \euro750\,000 for a firm dealing on its own account) and the \omterm{def:fm:the-proprietary-market-making-firm:ownfunds}{K-factor requirement} (article 15). For a dealer
the K-factors include K-DTF, 0.1\% of the daily trading flow in cash trades and 0.01\% in derivatives, the flow being the average of daily
absolute buys and sells over six months that end three months before the calculation (article 33), and either K-NPR, net position risk, or
K-CMG, the margin posted to a clearing member (article 21). The United Kingdom's MIFIDPRU rules follow the same design.
\textbf{United States.} SEC Rule 15c3-1: a broker-dealer's net capital is its net worth adjusted by, among other things, haircuts on its securities
positions; aggregate indebtedness may not exceed 1\,500\% of net capital, or, under the alternative standard, net capital must be at least the
greater of \$250\,000 and 2\% of aggregate debit items; a dealer keeps at least \$100\,000.
\end{dated}

\begin{definition}[Net capital rule]\label{def:fm:the-proprietary-market-making-firm:netcap}
The \emph{net capital rule}\index{net capital rule} is the United States requirement that a broker-dealer keep its net capital, its net worth less
illiquid assets and less haircuts on its positions, above a floor set by its business and by its indebtedness or its customers' debits.
\end{definition}

The European design ties a small firm's regulatory capital to its cost base, not to its risk. The chapter's firm has fixed overheads of \$95
million a year (people \$60 million, technology \$35 million) and trades \$10 billion a day of cash securities: its \omterm{def:fm:the-proprietary-market-making-firm:ownfunds}{fixed overheads requirement}
is \$23.75 million, its K-DTF \$10 million, its permanent minimum \$0.75 million. The binding requirement is the fixed overheads one
(\cref{fig:fm:the-proprietary-market-making-firm:capital}): a firm that hires or buys servers raises its regulatory capital before it has traded a
share more.

\begin{omfigure}
\begin{tikzpicture}
\begin{axis}[width=10.5cm,height=5.2cm,xbar,bar width=10pt,xmode=log,log basis x=10,log origin=infty,xmin=0.3,xmax=600,
  ytick={0,1,2,3,4},yticklabels={permanent minimum,fixed overheads (a quarter),K-DTF cash trades,K-DTF derivatives,prime brokers' requirement},
  y dir=reverse,ymin=-0.6,ymax=4.6,xlabel={\small \$ million (log scale)},tick label style={font=\footnotesize},table/col sep=comma,
  nodes near coords,nodes near coords style={font=\scriptsize,anchor=west},point meta=rawx,grid=major,grid style={gray!20}]
\addplot[fill=omDef!55,draw=omDef] table[y=k,x=amount]{figdata/desk/02-the-proprietary-market-making-firm/capital.csv};
\end{axis}
\end{tikzpicture}

\omcaption{The chapter's firm (fixed overheads \$95 million a year, \$10 billion a day of cash trades): the three parts of the European \omterm{def:fm:the-proprietary-market-making-firm:ownfunds}{own funds}
requirement, K-DTF if the same flow were derivatives, and the trading capital its prime brokers require. The regulatory requirement is the largest
of the first three; the prime brokers' is eight times it. Coefficients: Regulation (EU) 2019/2033; firm parameters illustrative. Data:
\texttt{fm\_partner.regulatory}.}\label{fig:fm:the-proprietary-market-making-firm:capital}
\end{omfigure}

The requirement that binds a market maker's growth is usually not the regulator's. Its prime brokers and clearing firms set their own
requirements, by internal haircut and margin models, and a market maker must post capital with them before it can hold positions; the listed
market maker of the hook calls this its trading capital and plans it daily. In the chapter's model the prime brokers require capital of 2\% of
the firm's daily volume: \$200 million supports \$10 billion a day, eight times the regulatory requirement.

\begin{dated}{2026-09}{One firm's trading capital}\label{dat:fm:the-proprietary-market-making-firm:flow}
Flow Traders Ltd.\ (Annual Report 2025): in July 2024 the board announced a Trading Capital Expansion Plan and suspended regular dividend
payments; no interim dividend was paid for 2024 or 2025. Net trading capital (net liquidity at clearing and prime brokers plus cash at bank)
was \euro1\,043.6 million at 31 December 2025, against \euro774.9 million a year earlier. In 2025 the firm took a \$200 million term loan and a \$75
million revolving facility. Its prime brokers ``require the Company to maintain certain minimum capital levels'', computed with internal haircut
and margin-based models. Net profit for 2025: \euro133.6 million; 635 full-time equivalents at the year end.
\end{dated}

\begin{method}[Sizing a trading firm's capital]\label{met:fm:the-proprietary-market-making-firm:size}
\begin{enumerate}
\item Compute the regulatory requirement: the highest of the permanent minimum, a quarter of last year's fixed overheads and the K-factors of the
planned flow and positions (\cref{lst:fm:the-proprietary-market-making-firm:req}).
\item Compute each prime broker's and clearing firm's requirement for the planned positions, on their models.
\item Add a buffer for the worst loss the firm plans to survive without cutting activity (chapters 7 and 8).
\item The firm needs the larger of the regulatory requirement and the sum of the posting requirements, plus the buffer; everything above is idle.
\end{enumerate}
\end{method}

\omcode{code/firm/partnership/firm_partnership.py}{78}{86}{The own funds requirement: the highest of its three parts.}\label{lst:fm:the-proprietary-market-making-firm:req}

\section{Reinvestment and payout}

The model of \texttt{firm.partnership} (\cref{lst:fm:the-proprietary-market-making-firm:sim}) runs a firm year by year: trading capital sets volume up
to what the market can give it; capture per unit traded depends on its technology; profit after fixed costs, technology spending and variable pay
is split into payout and retention. The chapter's firm starts with \$200 million, captures one basis point at technology parity, pays 25\% of its
result as variable pay, and its markets can give it at most \$20 billion a day, which \$400 million of capital supports. In its first year it trades
\$10 billion a day, earns \$250 million, spends \$35 million on technology and \$60 million on everything else, and makes \$116.25 million.

\omcode{code/firm/partnership/firm_partnership.py}{126}{143}{One year of the partnership: volume from capital, capture from technology, profit split
into payout and retention.}\label{lst:fm:the-proprietary-market-making-firm:sim}

\begin{proposition}[Retain until the market is full]\label{prop:fm:the-proprietary-market-making-firm:retain}
Suppose volume is $\min(K/m,\bar V)$ for capital $K$, a margin $m$ per unit of daily volume and a market of size $\bar V$, and profit per unit of
volume exceeds the partners' discount rate times $m$. Then every dollar retained while $K<m\bar V$ earns more than it would earn outside the firm,
and every dollar retained once $K\ge m\bar V$ earns nothing. Among policies that depend only on capital, retaining everything until $K$ reaches
$m\bar V$ and paying out everything afterwards maximises the discounted payout.
\end{proposition}

\begin{proof}[Partial proof]
Below capacity a retained dollar adds $1/m$ to daily volume and so a return of $\pi/m$ a year, $\pi$ the profit per unit of daily volume per
year; if $\pi/m$ exceeds the discount rate, deferring the payout raises its present value. Above capacity volume does not move and a retained
dollar returns nothing, while a paid-out dollar is worth its face value now. The policy that fills capacity as fast as possible and then pays
everything out is therefore optimal among capital-dependent policies; the model's fixed and variable costs do not change the argument because
they do not depend on retention. A full proof over all policies is a dynamic programme (One Quant Book 4, chapter 9).
\end{proof}

At the chapter's parameters the best fixed \omterm{def:fm:the-proprietary-market-making-firm:retention}{retention ratio} is 20\%, worth \$939 million of discounted payout over ten years at 10\%
(\cref{fig:fm:the-proprietary-market-making-firm:retention}); at 20\% the firm fills its market in year 6, at 50\% in year 3. Retaining everything
until capital reaches \$400 million and paying out everything afterwards is worth \$1\,210 million: the partners give up two years of payout and
gain \$271 million of value.

\begin{omfigure}
\begin{tikzpicture}
\begin{axis}[width=10.5cm,height=5.2cm,xlabel={\small retention ratio (\%)},ylabel={\small discounted payout (\$ million)},xmin=0,xmax=100,
  ymin=0,ymax=1350,tick label style={font=\footnotesize},grid=major,grid style={gray!20},table/col sep=comma,scaled ticks=false]
\addplot[omDef,thick,mark=*,mark size=1.2pt] table[x=retention,y=value]{figdata/desk/02-the-proprietary-market-making-firm/retention.csv};
\draw[omThm,dashed,thick] (axis cs:0,1209.9) -- (axis cs:100,1209.9);
\node[font=\scriptsize,omThm,anchor=south east,fill=white,inner sep=1pt] at (axis cs:98,1215) {retain until full, then pay out all};
\end{axis}
\end{tikzpicture}

\omcaption{Ten years of payouts discounted at 10\%, against a fixed \omterm{def:fm:the-proprietary-market-making-firm:retention}{retention ratio} (points), and for the policy of retaining everything until the
firm's capital fills its market and paying out everything afterwards (dashed line). Model: \texttt{fm\_partner.retention\_curve}, illustrative
parameters.}\label{fig:fm:the-proprietary-market-making-firm:retention}
\end{omfigure}

A departure is a retention decision the firm did not choose. If the member with 30 points leaves at the end of year 5 and is repaid over three
years, the firm (retaining 20\%) has the capital to fill its market by the end of year 8 instead of year 6. Staggered repayment is what makes the departure survivable; a
members' agreement that let a senior partner withdraw everything at once would give that partner the power to stop the firm growing.

\section{The technology treadmill}

A market maker's capture depends on how fast it is compared with the others (One Quant Book 11, chapter 9): who reaches a stale quote first,
whose quote is picked off. Speed is relative, so standing still is falling behind.

\begin{definition}[Technology treadmill]\label{def:fm:the-proprietary-market-making-firm:treadmill}
The \emph{technology treadmill}\index{technology treadmill} is the condition of a trading firm whose revenue depends on its technology relative to
its competitors': because they keep investing, it must keep spending at their rate to hold its capture constant, and spending less lowers its
capture even if its own systems do not change.
\end{definition}

\begin{proposition}[The spending that holds parity]\label{prop:fm:the-proprietary-market-making-firm:parity}
Let competitors' technology grow at a rate $g$ a year and the firm's technology stock $T_t$ depreciate at a rate $d$, so that
$T_{t+1}=(1-d)T_t+s_t$ for spending $s_t$. The firm's relative technology $T_t/F_t$, with $F_{t+1}=(1+g)F_t$, stays constant if and only if
$s_t=(g+d)\,T_t$, and then its spending grows at the rate $g$.
\end{proposition}

\begin{proof}
$T_{t+1}/F_{t+1}=T_t/F_t$ iff $(1-d)T_t+s_t=(1+g)T_t$, i.e.\ $s_t=(g+d)T_t$; then $T_{t+1}=(1+g)T_t$ and $s_{t+1}=(1+g)s_t$.
\end{proof}

At $g=15\%$ and $d=20\%$ the firm must spend 35\% of its technology stock every year, \$35 million in year 1 and \$123 million in year 10. If it
spends 25\%, its relative technology, and with it its capture, falls to 0.44 of parity by year 10
(\cref{fig:fm:the-proprietary-market-making-firm:treadmill}). With a market that does not grow, the treadmill eventually wins either way: the
parity spender's profit reaches zero in year 20, when spending that grows at 15\% a year catches a revenue that cannot grow; the lean spender's in
year 18, from falling capture (\cref{fig:fm:the-proprietary-market-making-firm:profit}). With a market growing 10\% a year the parity spender stays
profitable until year 56.

\begin{omfigure}
\begin{tikzpicture}
\begin{axis}[width=10.5cm,height=5cm,xlabel={\small year},ylabel={\small technology relative to competitors},xmin=1,xmax=25,ymin=0,ymax=1.1,
  tick label style={font=\footnotesize},grid=major,grid style={gray!20},table/col sep=comma,
  legend style={font=\scriptsize,at={(0.5,-0.3)},anchor=north,/tikz/every even column/.append style={column sep=8pt}},legend columns=2]
\addplot[omDef,thick] table[x=year,y=rel_steady]{figdata/desk/02-the-proprietary-market-making-firm/treadmill.csv};
\addplot[omThm,thick,dashed] table[x=year,y=rel_lean]{figdata/desk/02-the-proprietary-market-making-firm/treadmill.csv};
\legend{spending 35\% of the stock (parity),spending 25\%}
\end{axis}
\end{tikzpicture}

\omcaption{The treadmill: competitors' technology grows 15\% a year and the firm's depreciates 20\%. Spending 35\% of its stock holds parity;
spending 25\% lets relative technology, and capture with it, decay. Model: \texttt{fm\_partner.treadmill}.}\label{fig:fm:the-proprietary-market-making-firm:treadmill}
\end{omfigure}

\begin{omfigure}
\begin{tikzpicture}
\begin{axis}[width=10.5cm,height=5cm,xlabel={\small year},ylabel={\small profit (\$ million)},xmin=1,xmax=22,ymin=-350,ymax=300,
  tick label style={font=\footnotesize},grid=major,grid style={gray!20},table/col sep=comma,restrict x to domain=1:22,
  legend style={font=\scriptsize,at={(0.5,-0.3)},anchor=north,/tikz/every even column/.append style={column sep=8pt}},legend columns=2]
\addplot[omDef,thick] table[x=year,y=profit_steady]{figdata/desk/02-the-proprietary-market-making-firm/treadmill.csv};
\addplot[omThm,thick,dashed] table[x=year,y=profit_lean]{figdata/desk/02-the-proprietary-market-making-firm/treadmill.csv};
\draw[gray] (axis cs:1,0) -- (axis cs:22,0);
\legend{spending 35\% of the stock (parity),spending 25\%}
\end{axis}
\end{tikzpicture}

\omcaption{Profit of the two firms of \cref{fig:fm:the-proprietary-market-making-firm:treadmill} in a market that does not grow (retention 20\%).
The lean firm earns more only in the first year; both reach zero, in years 18 and 20. Model: \texttt{fm\_partner.treadmill}.}\label{fig:fm:the-proprietary-market-making-firm:profit}
\end{omfigure}

\begin{remark}[What the filings show, and do not]\label{rem:fm:the-proprietary-market-making-firm:filings}
The listed electronic market maker of chapter 1 reported communication and data costs of \$209.4 million in 2019 and \$249.2 million in 2025, up
19\% while its \omterm{def:fm:the-economics-of-a-trading-firm:ntr}{net trading revenue} swung between \$1.0 billion and \$2.3 billion; its costs nearly tripled between 2016 and 2019, but through two
acquisitions, not through the treadmill alone. A filing mixes the treadmill with growth, acquisitions and inflation. The treadmill is a mechanism
the published research supports: Budish, Cramton and Shim describe a continuing arms race for speed in which improvements are competed away, and
Baron, Brogaard, Hagströmer and Kirilenko found that firms which improved their relative speed earned more (One Quant Book 11, chapter 1).
\end{remark}

\section{Tutorial: ten years of a partnership}\label{tut:fm:the-proprietary-market-making-firm:tut}

\begin{tutorial}
\textbf{Goal.} Choose how much of each year's profit a trading partnership should keep, with its capital rules and the treadmill in the model.
\textbf{End state:} \cref{fig:fm:the-proprietary-market-making-firm:retention,fig:fm:the-proprietary-market-making-firm:profit}.
\begin{tutsteps}
\item \textbf{The requirement.} \texttt{fm\_partner.regulatory()} computes year 1's \omterm{def:fm:the-proprietary-market-making-firm:ownfunds}{own funds} requirement from fixed overheads (\$95 million) and
daily trading flow (\$10 billion): \$23.75 million, set by fixed overheads.
\item \textbf{The allocation.} \texttt{fp.Partnership} with members of 40, 30, 20 and 10 points allocates \$116.25 million
(\cref{ex:fm:the-proprietary-market-making-firm:alloc}); \texttt{fm\_partner.accounts()} runs ten years with a departure in year 5.
\item \textbf{Retention.} \texttt{fm\_partner.retention\_curve()} simulates ten years for \omterm{def:fm:the-proprietary-market-making-firm:retention}{retention ratios} from 0 to 100\% in steps of 5\% and
discounts the payouts at 10\%; \texttt{target\_value()} runs the fill-then-pay policy.
\item \textbf{The treadmill.} \texttt{fm\_partner.treadmill()} runs 25 years at 35\% and at 25\% of the technology stock.
\end{tutsteps}
\end{tutorial}

\textbf{What to change next.} Let the market grow (exercise 7); give the prime brokers a stricter margin in a volatile year and watch the
target capital move.

\section{Build: the partnership model}\label{bld:fm:the-proprietary-market-making-firm:partnership}

\begin{build}
\bfield{Purpose} The owners' side of a proprietary firm: who owns the capital, what the rules require of it, and what keeping it or paying it
out does to the firm.
\bfield{Interface} \texttt{firm.partnership}: \texttt{Member}, \texttt{Partnership} with \texttt{allocate(profit, payout\_ratio)},
\texttt{depart(name, years)}, \texttt{pay\_departures()}, \texttt{total\_capital}; \texttt{own\_funds\_requirement(fixed\_overheads, dtf\_cash,
dtf\_deriv, npr, \dots)}; \texttt{steady\_spend\_share(g, d)}, \texttt{capture(rel, c0, eta)}; \texttt{Params}, \texttt{simulate(params,
retention)} with a ratio or a policy, \texttt{target\_capital(target)}, \texttt{discounted\_payout(sim, rate, terminal\_multiple)}.
\bfield{Rules} A loss is allocated by points and reduces capital; nothing is drawn from a loss. A departing member's capital leaves in equal
instalments. Volume never exceeds capital over margin nor the market's size. Coefficients are inputs.
\bfield{Acceptance tests} \texttt{code/firm/partnership/tests/}: allocations by points, losses and staggered repayment; the requirement picks the
highest part; parity spending holds relative technology at one and lower spending decays it; capital changes by profit less payout; the
target-capital policy pays nothing until the target is reached.
\bfield{Stretch} A stochastic capture with losing years and a stop on trading when capital falls below the requirement; a term loan that adds
trading capital at an interest cost.
\end{build}

\begin{omsources}
\item Regulation (EU) 2019/2033 on the prudential requirements of investment firms (articles 11, 13--15, 21, 33) and Directive (EU) 2019/2034,
article 9.
\item US Code of Federal Regulations, 17 CFR 240.15c3-1 (net capital rule).
\item Limited Liability Partnerships Act 2000 (United Kingdom), section 1.
\item Flow Traders Ltd., Annual Report 2025.
\item E.~Budish, P.~Cramton and J.~Shim, ``The high-frequency trading arms race: frequent batch auctions as a market design response'',
\emph{Quarterly Journal of Economics} 130(4), 2015.
\end{omsources}

\section{Exercises}

\begin{exercise}[$\star$]\label{exo:fm:the-proprietary-market-making-firm:1}
A dealer on its own account had fixed overheads of \$95 million last year, trades \$10 billion a day of cash securities and holds no overnight
positions. Compute the three parts of its European \omterm{def:fm:the-proprietary-market-making-firm:ownfunds}{own funds} requirement and say which binds.
\end{exercise}

\begin{exercise}[$\star$]\label{exo:fm:the-proprietary-market-making-firm:2}
The same firm trades \$10 billion a day of derivatives notional instead. What is its K-DTF, and does the binding requirement change?
\end{exercise}

\begin{exercise}[$\star$]\label{exo:fm:the-proprietary-market-making-firm:3}
Allocate a profit of \$116.25 million to members of 40, 30, 20 and 10 points who draw 80\% of their allocations. What does each draw, and what
stays in the firm?
\end{exercise}

\begin{exercise}[$\star\star$]\label{exo:fm:the-proprietary-market-making-firm:4}
Competitors' technology grows 15\% a year and the firm's depreciates 20\%. What share of its stock must it spend to hold parity, and how much is
that in year 10 if it spends \$35 million in year 1?
\end{exercise}

\begin{exercise}[$\star\star$]\label{exo:fm:the-proprietary-market-making-firm:5}
The market is full at \$20 billion a day and one basis point of capture, fixed costs are \$60 million and parity spending starts at \$35 million
and grows 15\% a year. In which year does spending make profit before variable pay zero? Derive it in closed form.
\end{exercise}

\begin{exercise}[$\star\star$]\label{exo:fm:the-proprietary-market-making-firm:6}
In \cref{fig:fm:the-proprietary-market-making-firm:retention}, why does the value fall for \omterm{def:fm:the-proprietary-market-making-firm:retention}{retention ratios} above 20\%, and why is the fixed-ratio
curve always below the dashed line?
\end{exercise}

\begin{exercise}[$\star\star\star$]\label{exo:fm:the-proprietary-market-making-firm:7}
\emph{Coding.} Rerun the parity spender for 60 years with a market growing 10\% a year (\texttt{fm\_partner.growing\_market}). In which year does
its profit first fall to zero, and why does it fall at all?
\end{exercise}

\begin{exercise}[$\star\star\star$]\label{exo:fm:the-proprietary-market-making-firm:8}
\emph{Find the flaw.} ``We cut technology spending from 35\% to 25\% of our stock last year and profit rose from \$116 million to \$124 million.
Cut it again.''
\end{exercise}

\section{Problem: Keep It or Pay It Out}

\begin{problem}[{Weekend problem --- keep it or pay it out}]\label{pb:fm:the-proprietary-market-making-firm:1}
Four partners run a market-making firm with \$200 million of capital. They must decide how much of each year's profit to keep, and how much to
spend on technology.

\medskip\noindent\textbf{Part I --- The rules.}
\begin{enumerate}
\item Define \omterm{def:fm:the-proprietary-market-making-firm:ownfunds}{own funds} and state the European \omterm{def:fm:the-proprietary-market-making-firm:ownfunds}{own funds} requirement.
\item Compute the firm's \omterm{def:fm:the-proprietary-market-making-firm:ownfunds}{fixed overheads requirement} in year 1.
\item Compute its K-DTF on \$10 billion a day of cash trades, and on the same notional of derivatives.
\item Which requirement binds, and how does it compare with the \$200 million its prime brokers require?
\item State the US \omterm{def:fm:the-proprietary-market-making-firm:netcap}{net capital rule}'s mechanism in one sentence.
\end{enumerate}

\medskip\noindent\textbf{Part II --- The partnership.}
\begin{enumerate}[resume]
\item Define a \omterm{def:fm:the-proprietary-market-making-firm:llp}{limited liability partnership} and \omterm{def:fm:the-proprietary-market-making-firm:llp}{members' capital}.
\item Compute year 1's volume, revenue, technology spending and profit.
\item Allocate year 1's profit to the 40-point member at 80\% payout: her draw and what she leaves in.
\item The 30-point member leaves at the end of year 5, repaid over three years. In which year does the firm fill its market (retention 20\%),
against which year without the departure?
\end{enumerate}

\medskip\noindent\textbf{Part III --- Retention.}
\begin{enumerate}[resume]
\item Define the \omterm{def:fm:the-proprietary-market-making-firm:retention}{retention ratio}.
\item Give the ten-year discounted payout at \omterm{def:fm:the-proprietary-market-making-firm:retention}{retention ratios} of 0, 20, 50 and 80\%.
\item Give the best fixed \omterm{def:fm:the-proprietary-market-making-firm:retention}{retention ratio} and its value.
\item Give the value of retaining everything until the market is full, and explain the gain with \cref{prop:fm:the-proprietary-market-making-firm:retain}.
\item In which year is the market full at \omterm{def:fm:the-proprietary-market-making-firm:retention}{retention ratios} of 20, 50 and 80\%?
\end{enumerate}

\medskip\noindent\textbf{Part IV --- The treadmill.}
\begin{enumerate}[resume]
\item Define the \omterm{def:fm:the-proprietary-market-making-firm:treadmill}{technology treadmill}.
\item Derive the spending share that holds parity.
\item What is the lean spender's capture after ten years, as a fraction of parity?
\item In which year does each firm's profit reach zero in a market that does not grow?
\item State the \emph{named result}: the parity spending share and its year-10 amount, the best retention policy and its value against the best
fixed ratio.
\item In two sentences, what should the partners do first?
\end{enumerate}
\end{problem}

\section{Interview questions}

\begin{interviewq}[$\star$ \iqroles{trader, risk}]\label{iq:fm:the-proprietary-market-making-firm:1}
Under the European rules, why can a small proprietary firm's regulatory capital rise when it hires, even if it trades no more?
\end{interviewq}

\begin{interviewq}[$\star$ \iqroles{trader}]\label{iq:fm:the-proprietary-market-making-firm:2}
Which capital requirement usually binds a market maker's growth: the regulator's or its prime brokers'? Why?
\end{interviewq}

\begin{interviewq}[$\star\star$ \iqroles{trader, researcher}]\label{iq:fm:the-proprietary-market-making-firm:3}
Your firm can use \$400 million of trading capital and has \$250 million. Should the partners draw this year's profit?
\end{interviewq}

\begin{interviewq}[$\star\star$ \iqroles{developer}]\label{iq:fm:the-proprietary-market-making-firm:4}
Competitors' technology improves 10\% a year and yours depreciates 25\%. What share of your technology stock must you spend each year to keep up?
\end{interviewq}

\begin{interviewq}[$\star\star$ \iqroles{risk}]\label{iq:fm:the-proprietary-market-making-firm:5}
Why do members' agreements repay a departing partner's capital over several years?
\end{interviewq}

\begin{interviewq}[$\star\star\star$ \iqroles{trader, developer}]\label{iq:fm:the-proprietary-market-making-firm:6}
A firm cut its technology budget and its profit rose the next year. What would you measure before concluding the cut was right?
\end{interviewq}
